% --- Archon physics formalization source begin ---
% archon:physics
% archon:covers PhyXMiniProblems/problem_phyx_mini_0886.lean
% archon:source-report reports/phyx_mini/problem_phyx_mini_0886.source.json

\chapter{Physics problem phyx\_mini\_0886}
\label{ch:PhyXMiniProblems_problem_phyx_mini_0886}

\paragraph{Problem source.}
\#\# Physical scenario

The emf phasor in figure is shown at \textbackslash{}( t = 15\textbackslash{},\textbackslash{}text\{ms\} \textbackslash{}).

\#\# Question

What is the instantaneous value of the emf?

\#\# Answer choices

- A: A: \textbackslash{}( -6\textbackslash{}

- B: B: \textbackslash{}( -8\textbackslash{}

- C: C: \textbackslash{}( -10\textbackslash{}

- D: D: \textbackslash{}( -5\textbackslash{}

\#\# Figure caption (auxiliary; use the image as primary evidence)

The image depicts a phasor diagram. 

- It has a coordinate system with horizontal and vertical axes intersecting at the origin.
- A green arrow, representing the phasor, originates from the origin and points towards the top left quadrant.
- The arrow is labeled with \textbackslash{}(\textbackslash{}mathcal\{E\}\_0 = 12 \textbackslash{}, \textbackslash{}text\{V\}\textbackslash{}) at its beginning, indicating the magnitude of the phasor.
- The angle between the arrow and the positive horizontal axis is marked as \textbackslash{}(30\textasciicircum{}\textbackslash{}circ\textbackslash{}).
- There is a text annotation, "Phasor at \textbackslash{}(t = 15 \textbackslash{}, \textbackslash{}text\{ms\}\textbackslash{})," situated above the arrow, indicating the time at which this phasor is represented.

\#\# Dataset metadata

Electromagnetic Waves; Physical Model Grounding Reasoning, Numerical Reasoning

\paragraph{Recorded answer/context.}
C: C: \textbackslash{}( -10\textbackslash{}

\paragraph{Figure/image path.}
/root/proposal\_for\_physic/science-mango/phyx\_mini\_run/phyx\_data/test\_image/886.png

\paragraph{Formalization target.}
create a compiling Lean file with sorry bodies at `PhyXMiniProblems/problem\_phyx\_mini\_0886.lean`. The Lean declarations must preserve the physical quantities, dimensions or dimensional roles, figure labels, governing-law hypotheses, and final relation expressed by this problem.
Use Mathlib/Physlib names found through LeanExplore where available. If a physics API is missing, introduce faithful local abstractions rather than scalar placeholder aliases.

\begin{theorem}[Physics formalization target]
\label{thm:physics:phyx_mini_0886:target}
\lean{PhyXMiniProblems.ProblemPhyXMini0886.problem_phyx_mini_0886}
\uses{def:physics:phyx-mini-0886:phyxminiproblems-problemphyxmini0886-signedemfinvolts, def:physics:phyx-mini-0886:phyxminiproblems-problemphyxmini0886-emfphasorsetup, def:physics:phyx-mini-0886:phyxminiproblems-problemphyxmini0886-matchessuppliedemfphasorfigure, def:physics:phyx-mini-0886:phyxminiproblems-problemphyxmini0886-satisfiescosinereferencephasorlaw, def:physics:phyx-mini-0886:phyxminiproblems-problemphyxmini0886-recordeddatasetanswer, def:physics:phyx-mini-0886:phyxminiproblems-problemphyxmini0886-isuniquenearestdisplayedchoice}
The assigned autoformalize agent should translate this physics problem into Lean declarations in the covered file, with theorem and lemma proof bodies written as `by sorry`.
\end{theorem}
\begin{proof}
This is an autoformalization task, not a proof task. Produce statements that can later be proved by the physics prover without weakening the physical model.
\end{proof}

% --- Archon declaration topology begin ---
\section{Declaration topology}
\begin{definition}[emf Dimension]
  \label{def:physics:phyx-mini-0886:phyxminiproblems-problemphyxmini0886-emfdimension}
  \lean{PhyXMiniProblems.ProblemPhyXMini0886.emfDimension}
  The physical dimension M L² T⁻² C⁻¹ shared by EMF and voltage
\end{definition}
\begin{proof}
  This is the declared model definition of emf Dimension.
\end{proof}
\begin{definition}[Emf Magnitude Quantity]
  \label{def:physics:phyx-mini-0886:phyxminiproblems-problemphyxmini0886-emfmagnitudequantity}
  \lean{PhyXMiniProblems.ProblemPhyXMini0886.EmfMagnitudeQuantity}
  \uses{def:physics:phyx-mini-0886:phyxminiproblems-problemphyxmini0886-emfdimension}
  A nonnegative, unit-independent peak EMF magnitude
\end{definition}
\begin{proof}
  This is the declared model definition of Emf Magnitude Quantity.
\end{proof}
\begin{definition}[Signed Emf Quantity]
  \label{def:physics:phyx-mini-0886:phyxminiproblems-problemphyxmini0886-signedemfquantity}
  \lean{PhyXMiniProblems.ProblemPhyXMini0886.SignedEmfQuantity}
  \uses{def:physics:phyx-mini-0886:phyxminiproblems-problemphyxmini0886-emfdimension}
  A signed, unit-independent instantaneous EMF
\end{definition}
\begin{proof}
  This is the declared model definition of Signed Emf Quantity.
\end{proof}
\begin{definition}[Time Quantity]
  \label{def:physics:phyx-mini-0886:phyxminiproblems-problemphyxmini0886-timequantity}
  \lean{PhyXMiniProblems.ProblemPhyXMini0886.TimeQuantity}
  A nonnegative, unit-independent time
\end{definition}
\begin{proof}
  This is the declared model definition of Time Quantity.
\end{proof}
\begin{definition}[emf Magnitude In Volts]
  \label{def:physics:phyx-mini-0886:phyxminiproblems-problemphyxmini0886-emfmagnitudeinvolts}
  \lean{PhyXMiniProblems.ProblemPhyXMini0886.emfMagnitudeInVolts}
  \uses{def:physics:phyx-mini-0886:phyxminiproblems-problemphyxmini0886-emfmagnitudequantity}
  Coherent-SI readout of a peak EMF magnitude in volts
\end{definition}
\begin{proof}
  This is the declared model definition of emf Magnitude In Volts.
\end{proof}
\begin{definition}[signed Emf In Volts]
  \label{def:physics:phyx-mini-0886:phyxminiproblems-problemphyxmini0886-signedemfinvolts}
  \lean{PhyXMiniProblems.ProblemPhyXMini0886.signedEmfInVolts}
  \uses{def:physics:phyx-mini-0886:phyxminiproblems-problemphyxmini0886-signedemfquantity}
  Coherent-SI readout of a signed instantaneous EMF in volts
\end{definition}
\begin{proof}
  This is the declared model definition of signed Emf In Volts.
\end{proof}
\begin{definition}[time In Milliseconds]
  \label{def:physics:phyx-mini-0886:phyxminiproblems-problemphyxmini0886-timeinmilliseconds}
  \lean{PhyXMiniProblems.ProblemPhyXMini0886.timeInMilliseconds}
  \uses{def:physics:phyx-mini-0886:phyxminiproblems-problemphyxmini0886-timequantity}
  Millisecond readout used by the time label in the supplied image
\end{definition}
\begin{proof}
  This is the declared model definition of time In Milliseconds.
\end{proof}
\begin{definition}[Cartesian Quadrant]
  \label{def:physics:phyx-mini-0886:phyxminiproblems-problemphyxmini0886-cartesianquadrant}
  \lean{PhyXMiniProblems.ProblemPhyXMini0886.CartesianQuadrant}
  The four open quadrants determined by the two axes in the phasor image
\end{definition}
\begin{proof}
  This is the declared model definition of Cartesian Quadrant.
\end{proof}
\begin{definition}[Emf Phasor Figure]
  \label{def:physics:phyx-mini-0886:phyxminiproblems-problemphyxmini0886-emfphasorfigure}
  \lean{PhyXMiniProblems.ProblemPhyXMini0886.EmfPhasorFigure}
  \uses{def:physics:phyx-mini-0886:phyxminiproblems-problemphyxmini0886-cartesianquadrant}
  Literal presentation data transcribed from image 886.png
\end{definition}
\begin{proof}
  This is the declared model definition of Emf Phasor Figure.
\end{proof}
\begin{definition}[Emf Phasor Setup]
  \label{def:physics:phyx-mini-0886:phyxminiproblems-problemphyxmini0886-emfphasorsetup}
  \lean{PhyXMiniProblems.ProblemPhyXMini0886.EmfPhasorSetup}
  \uses{def:physics:phyx-mini-0886:phyxminiproblems-problemphyxmini0886-emfmagnitudequantity, def:physics:phyx-mini-0886:phyxminiproblems-problemphyxmini0886-signedemfquantity, def:physics:phyx-mini-0886:phyxminiproblems-problemphyxmini0886-timequantity, def:physics:phyx-mini-0886:phyxminiproblems-problemphyxmini0886-emfphasorfigure}
  Independent physical quantities and observables for the phasor experiment. In particular, instantaneousEmf is not defined from the phasor projection, the recorded answer, or any displayed numerical choice
\end{definition}
\begin{proof}
  This is the declared model definition of Emf Phasor Setup.
\end{proof}
\begin{definition}[Matches Supplied Emf Phasor Figure]
  \label{def:physics:phyx-mini-0886:phyxminiproblems-problemphyxmini0886-matchessuppliedemfphasorfigure}
  \lean{PhyXMiniProblems.ProblemPhyXMini0886.MatchesSuppliedEmfPhasorFigure}
  \uses{def:physics:phyx-mini-0886:phyxminiproblems-problemphyxmini0886-emfmagnitudeinvolts, def:physics:phyx-mini-0886:phyxminiproblems-problemphyxmini0886-timeinmilliseconds, def:physics:phyx-mini-0886:phyxminiproblems-problemphyxmini0886-emfphasorsetup}
  Primary-image evidence. The raster places the marked 30° between the arrow and the negative horizontal ray, despite the auxiliary caption's wording about the positive horizontal axis. The final field converts that pictorial relation into the cosine-reference phase at the displayed time
\end{definition}
\begin{proof}
  This is the declared model definition of Matches Supplied Emf Phasor Figure.
\end{proof}
\begin{definition}[Satisfies Cosine Reference Phasor Law]
  \label{def:physics:phyx-mini-0886:phyxminiproblems-problemphyxmini0886-satisfiescosinereferencephasorlaw}
  \lean{PhyXMiniProblems.ProblemPhyXMini0886.SatisfiesCosineReferencePhasorLaw}
  \uses{def:physics:phyx-mini-0886:phyxminiproblems-problemphyxmini0886-emfmagnitudeinvolts, def:physics:phyx-mini-0886:phyxminiproblems-problemphyxmini0886-signedemfinvolts, def:physics:phyx-mini-0886:phyxminiproblems-problemphyxmini0886-emfphasorsetup}
  For a cosine-reference phasor, the signed instantaneous EMF is its horizontal projection. This law is uniform in time and contains neither the evaluated value -6 * sqrt 3 nor a displayed answer choice
\end{definition}
\begin{proof}
  This is the declared model definition of Satisfies Cosine Reference Phasor Law.
\end{proof}
\begin{definition}[Answer Choice]
  \label{def:physics:phyx-mini-0886:phyxminiproblems-problemphyxmini0886-answerchoice}
  \lean{PhyXMiniProblems.ProblemPhyXMini0886.AnswerChoice}
  Labels of the four displayed answer choices
\end{definition}
\begin{proof}
  This is the declared model definition of Answer Choice.
\end{proof}
\begin{definition}[emf In Volts]
  \label{def:physics:phyx-mini-0886:phyxminiproblems-problemphyxmini0886-answerchoice-emfinvolts}
  \lean{PhyXMiniProblems.ProblemPhyXMini0886.AnswerChoice.emfInVolts}
  \uses{def:physics:phyx-mini-0886:phyxminiproblems-problemphyxmini0886-answerchoice}
  Voltage value printed beside each displayed answer label
\end{definition}
\begin{proof}
  This is the declared model definition of emf In Volts.
\end{proof}
\begin{definition}[recorded Dataset Answer]
  \label{def:physics:phyx-mini-0886:phyxminiproblems-problemphyxmini0886-recordeddatasetanswer}
  \lean{PhyXMiniProblems.ProblemPhyXMini0886.recordedDatasetAnswer}
  \uses{def:physics:phyx-mini-0886:phyxminiproblems-problemphyxmini0886-answerchoice}
  The answer label recorded in the supplied dataset metadata
\end{definition}
\begin{proof}
  This is the declared model definition of recorded Dataset Answer.
\end{proof}
\begin{definition}[Is Unique Nearest Displayed Choice]
  \label{def:physics:phyx-mini-0886:phyxminiproblems-problemphyxmini0886-isuniquenearestdisplayedchoice}
  \lean{PhyXMiniProblems.ProblemPhyXMini0886.IsUniqueNearestDisplayedChoice}
  \uses{def:physics:phyx-mini-0886:phyxminiproblems-problemphyxmini0886-answerchoice}
  A displayed choice is uniquely nearest to a computed scalar readout when it has strictly smaller absolute error than every differently labelled choice. This generic definition does not select C or prescribe the physical EMF
\end{definition}
\begin{proof}
  This is the declared model definition of Is Unique Nearest Displayed Choice.
\end{proof}
\begin{lemma}[phase At Observation Time eq five pi div six]
  \label{lem:physics:phyx-mini-0886:phyxminiproblems-problemphyxmini0886-phaseatobservationtime-eq-five-pi-div-six}
  \lean{PhyXMiniProblems.ProblemPhyXMini0886.phaseAtObservationTime_eq_five_pi_div_six}
  \uses{def:physics:phyx-mini-0886:phyxminiproblems-problemphyxmini0886-emfphasorsetup, def:physics:phyx-mini-0886:phyxminiproblems-problemphyxmini0886-matchessuppliedemfphasorfigure}
  The second-quadrant geometry makes the phase at the labelled time 5π/6 modulo a full turn
\end{lemma}
\begin{proof}
  The conclusion follows from the governing relations and figure readouts named in the statement of phase At Observation Time eq five pi div six.
\end{proof}
% --- Archon declaration topology end ---
% --- Archon physics formalization source end ---
